\section{Protocol}
\label{sec:protocol}

\begin{figure}[!htbp]
\centering
\begin{tikzpicture}[
  font=\footnotesize, node distance=0.25cm,
  tier/.style={draw, rounded corners=2pt, align=left, inner sep=4pt, text width=7.35cm, minimum height=1.55cm},
  arr/.style={-{Stealth[length=1.8mm]}, thick}
]
\node[tier, fill=black!4] (a) {\textbf{Tier A --- reproduction (RQ1, H1)}\\ 30 fev-bench tasks; Google's published CSV;
  pipeline controls on Seasonal Naive and Chronos-2};
\node[tier, fill=blue!8, right=of a] (b) {\textbf{Tier B --- post-release real data (RQ2, family P, H5)}\\ load, price, weather,
  river, FX, FRED; windows 2026-09-01 to 10-20, after the 2026-08-31 release};
\node[tier, fill=blue!8, below=of a] (w) {\textbf{Tier W --- memorisation probe (H2)}\\ Wikipedia pageviews: same pages in 2023
  (W-seen) and 2026 (W-post); pages born after 2023 (W-new)};
\node[tier, fill=blue!8, below=of b] (p) {\textbf{Tier P --- prospective check}\\ two issues (2026-10-30, 11-02); forecasts
  sealed by SHA-256 before the outcomes existed};
\node[tier, fill=orange!12, below=of w] (c) {\textbf{Tier C --- controlled synthetic data (RQ3, H3, H4)}\\ 38 cells $\times$ 200
  replications: failure modes, multivariate structure, replication of a published simulation};
\node[tier, fill=orange!12, below=of p] (abl) {\textbf{Ablations (RQ6)}\\ eight configurations, each switching one fixed step or
  option, on Tier C and a 20\% sample of Tier B};
\node[draw, rounded corners=2pt, fill=black!2, align=center, text width=15.3cm, inner sep=4pt, below=0.45cm of $(c.south)!0.5!(abl.south)$]
  (out) {\textbf{Same scoring for every tier:} MASE and SQL relative to Seasonal Naive $\cdot$ paired log ratios $\cdot$ two-way
  cluster bootstrap $\cdot$ Holm within pre-declared families $\cdot$ protocol frozen by a git tag before any post-release data};
\end{tikzpicture}
\caption{Study design. Tiers B and P, and the W-post and W-new strata of Tier W, use data published after TimesFM~3's release (blue; W-seen uses 2023 windows); Tier A is the public benchmark the
model was evaluated on, and Tier C and the ablations use synthetic data whose generating process is known (orange).}
\label{fig:design}
\end{figure}

\begin{figure}[!t]
\centering
\begin{tikzpicture}[x=0.143cm, y=1cm, font=\footnotesize,
  lane/.style={font=\footnotesize\bfseries, text=vzink, anchor=east},
  ev/.style={circle, fill=#1, draw=white, line width=1pt, inner sep=2.1pt},
  seal/.style={diamond, fill=#1, draw=white, line width=1pt, inner sep=2.2pt},
  lab/.style={font=\scriptsize, text=vzink2, align=#1},
  band/.style={fill=#1, rounded corners=1.5pt}]
% day 0 = 2026-08-20
\foreach \y in {0,-0.8,-1.6,-2.4} \draw[vzgrid, line width=0.4pt] (0,\y) -- (89,\y);
\foreach \d/\m in {12/1 Sep, 42/1 Oct, 73/1 Nov} {
  \draw[vzaxis, line width=0.4pt] (\d,0.45) -- (\d,-3.05);
  \node[font=\scriptsize, text=vzmuted, anchor=south] at (\d,0.45) {\m};}
\node[lane] at (-0.8,0) {Model};
\node[lane] at (-0.8,-0.8) {Test data};
\node[lane] at (-0.8,-1.6) {Protocol};
\node[lane, align=right] at (-0.8,-2.6) {Prospective\\check};
% model
\node[ev=vzmuted] (w) at (4,0) {};
\node[ev=vzink] (r) at (11,0) {};
\node[lab=left, anchor=west] at (12,0) {\textbf{TimesFM~3 released, 31 Aug}\quad (grey: weights first uploaded, 24 Aug)};
% Tier B windows
\fill[band=vzblue] (12,-0.92) rectangle (61,-0.68);
\node[font=\scriptsize, text=white, anchor=center] at (36.5,-0.8) {Tier B windows, 1 Sep -- 20 Oct (six domains)};
% protocol
\node[ev=vzink] (f) at (69,-1.6) {};
\node[lab=right, anchor=east] at (68,-1.6) {protocol frozen (git tag), 28 Oct;\\post-release data downloaded that evening};
\node[ev=vzaqua] at (71,-1.6) {};
\node[lab=left, anchor=west] at (72,-1.6) {first public\\timestamp, 30 Oct};
% prospective issues
\node[seal=vzorange] at (71,-2.3) {};
\fill[band=vzorange!55] (73,-2.38) rectangle (80,-2.22);
\node[ev=vzorange] at (82,-2.3) {};
\node[lab=right, anchor=east] at (70,-2.3) {issue 1 sealed (SHA-256), 30 Oct};
\node[lab=left, anchor=west] at (83,-2.3) {10 Nov};
\node[seal=vzorange] at (74,-2.9) {};
\fill[band=vzorange!55] (76,-2.98) rectangle (83,-2.82);
\node[circle, fill=white, draw=vzorange, line width=1pt, inner sep=2.1pt] at (85,-2.9) {}; % planned: switch to ev=vzorange once scored
\node[lab=right, anchor=east] at (73,-2.9) {issue 2 sealed, 2 Nov};
\node[lab=left, anchor=west] at (86,-2.9) {13 Nov};
% legend
\node[font=\scriptsize, text=vzink2, anchor=north west, align=left] at (0,-3.25)
  {\tikz[baseline=-0.6ex]{\pgfsetfillcolor{vzorange}\pgfsetstrokecolor{vzorange}\pgfsetplotmarksize{2.6pt}\pgfuseplotmark{diamond*}}~forecasts sealed \quad
   \tikz[baseline=-0.6ex]\fill[vzorange!55, rounded corners=1pt] (0,-0.08) rectangle (0.5,0.08);~outcome window\quad
   \tikz[baseline=-0.6ex]\fill[vzorange] (0,0) circle (2.1pt);~outcomes scored\quad
   \tikz[baseline=-0.6ex]\draw[vzorange, line width=1pt] (0,0) circle (2.1pt);~planned\\
   \tikz[baseline=-0.6ex]\fill[vzaqua] (0,0) circle (2.1pt);~external timestamp (RFC~3161 and OpenTimestamps)};
\end{tikzpicture}
\caption{Chronology, 2026. Every Tier B test window starts after TimesFM~3's release; the protocol was frozen by a git tag
before the post-release data were downloaded (an internal record; its first external timestamp is two days later); the two
prospective issues were sealed by hash before their outcomes existed and are scored after the outcome windows close.}
\label{fig:timeline}
\end{figure}


\paragraph{Freeze.} Research questions, hypotheses, series lists, metrics, contrasts, families and decision rules were written in a
protocol and frozen by the git tag \texttt{protocol-frozen-v1} (commit \texttt{9f15b38}, 2026-10-28 17:59 UTC) before any data of
the post-release test windows were downloaded; the download code refuses earlier dates. The freeze is internal: the hash of the
frozen commit received its first public timestamp on 2026-10-30, after the download, so the order of freeze and download rests
on the repository history, not on an external record (Figure~\ref{fig:timeline}). The protocol incorporates an independent
adversarial review of its draft and a pilot on pre-release data (one year earlier) and on a separate synthetic set; the pilot is not
part of the results. Two deviations occurred (Appendix~\ref{app:deviations}): the post-release data were downloaded on the evening of
the freeze day instead of the next day, and a software error voided one ablation, which was re-run. Interpretations of
ambiguous protocol items are listed with them. Results that the protocol did not specify are labelled \emph{exploratory} or
\emph{descriptive} wherever they appear.

\paragraph{Research questions.} RQ1: do the published per-task TimesFM~3 results, and the claim they support, reproduce with the
public code and weights? RQ2: does TimesFM~3's advantage hold on data that post-date its release, and is there a trace of
memorisation? RQ3: under which controlled conditions does zero-shot forecasting fail? RQ4: are the quantile forecasts calibrated?
RQ5: when do multivariate mode and covariates help? RQ6: how much do inference-time options and fixed components contribute? RQ7: is
zero-shot forecasting worth it against classical methods, accounting for compute?

\paragraph{Pre-declared hypotheses.} \textbf{H1} (descriptive): the advantage over a statistical ensemble is smaller on post-release
data than on the public benchmark. \textbf{H2}: on Wikipedia pageviews, TimesFM~3's advantage over Chronos-2 is larger on
2023 windows inside TimesFM's pageview corpus (W-seen) than on the same pages in 2026 (W-post). \textbf{H3}: disabling the linear
detrending lowers the error over steps 25--48 on a saturating logistic process. \textbf{H4}: multivariate mode lowers the error
when a known dependence exists (H4a lead-lag, H4b known future covariate) and is equivalent within $\pm$3\% when none exists (H4c
independent series, H4d spurious covariate). \textbf{H5}: on near-random-walk series (exchange rates, Treasury yields, oil prices)
TimesFM~3 does not improve on the random walk by more than 10\%.

\paragraph{Metrics.} The point metric is MASE of the median forecast and the probabilistic metric is the scaled quantile loss (SQL)
over the nine quantile levels, both as defined by fev-bench \citep{shchur2025fevbench}; scales are the in-sample seasonal
differences of each task's full context (season 24 for hourly, 7 for daily, 5 for business-daily data), the same for every model
whatever context it is given, so the scale cancels in a paired ratio. We also record RMSSE, WAPE, WQL and the
empirical coverage of the 80\% and 60\% central intervals.

\paragraph{Estimand and uncertainty.} A contrast between configurations A and B is the paired log ratio
$d = \log(\text{error}_A) - \log(\text{error}_B)$ on identical series and windows, averaged with equal weight within each domain
$\times$ horizon cell, then over horizons and over domains, and reported as $\exp(\bar d)$: a ratio below 1 favours A. Uncertainty
comes from a two-way cluster bootstrap with 10{,}000 replicates that resamples series (within domains) and ISO weeks of the first
forecast step jointly \citep{cameron2011multiway}: each replicate draws multinomial weights for series and for weeks and multiplies them
(a pigeonhole bootstrap, \citealp{owen2007pigeonhole}), so the dependence within a series and the shocks common to several series in the same week are reflected.
Dependence between different series in different but overlapping windows is not, and with eight weekly clusters temporal
inference is weak (Section~\ref{sec:robustness}). Synthetic data resample replications. Pairs where either error is zero or not
finite are dropped (the log ratio is undefined). Intervals are percentile 95\% intervals; $p$-values are two-sided, twice the smaller share of replicates on either side
of zero, with a floor of $1/10{,}000$ (seed 17); the series-only interval is reported as a sensitivity analysis.

\paragraph{Families and decisions.} Decisions use Holm-adjusted $p$-values \citep{holm1979} within a family at
$\alpha = 0.05$; intervals are unadjusted. \emph{Family P} (primary, ten tests) holds five contrasts on
post-release data, each on MASE and SQL: TimesFM~3 in univariate mode without covariates against the statistical ensemble (P1),
Chronos-2 (P2) and TimesFM~2.5 (P3), all without covariates, on the same 3{,}855 univariate evaluations of the four structured
domains (on hourly data the context lengths differ; Section~\ref{sec:data}); the multivariate against the univariate mode on the
884 evaluations of the price, FX and FRED group tasks, which have no covariates (P4a); and the benchmark flags against the package
defaults on the 3{,}855 univariate evaluations, both configurations with the covariates of load and river flow (P5). \emph{Family H} holds the primary test of each hypothesis H2--H5
(seven tests). Equivalence (H4c, H4d) is claimed only if the 90\% interval lies inside the $\pm$3\% margin (two one-sided tests
\citep{schuirmann1987tost}); H5 is a non-superiority test that is supported when the 95\% lower bound exceeds 0.90. A
difference whose adjusted $p$ exceeds 0.05 is reported as \emph{not resolved}, never as a tie. Secondary families (S1--S6) have their
own Holm correction.
